\chapter{Questions You Ask a Protocol}
\label{ch:questions}
\label{ch:ctl}

\index{CTL}
\index{temporal logic}
Some rules are visible in a protocol's shape: who decides, who is told, which
messages are calls. Others are not. ``Never charge a card after saying sold
out'' and ``every charge ends in a receipt'' are questions about every possible
future of a conversation. LeanFM asks them in CTL, computation tree logic
\citep{clarke1986automatic}, over the protocol's own state graph, and every answer
comes with a path: a witness when something is possible, a counterexample when
something fails.

\section{The operators}

A CTL operator pairs a path quantifier (\textbf{A}, on every path; \textbf{E}, on
some path) with a time operator (\textbf{F}, eventually; \textbf{G}, always;
\textbf{U}, until).

\begin{center}
\small
\begin{tabular}{@{}llp{.55\textwidth}@{}}
\toprule
Formula & Read as & Meaning \\
\midrule
\texttt{EF p} & possibly & some future reaches \texttt{p} \\
\texttt{AF p} & inevitably & every future reaches \texttt{p} \\
\texttt{AG p} & always & \texttt{p} holds now and in every future state \\
\texttt{AG not p}, \texttt{never p} & never & no future reaches \texttt{p} \\
\texttt{EG p} & can stay & some future keeps \texttt{p} forever \\
\texttt{A[p U q]}, \texttt{E[p U q]} & until & \texttt{p} holds until \texttt{q} does, on every path or some \\
\texttt{AW}, \texttt{EW} & unless & until, without requiring that \texttt{q} ever happens \\
\bottomrule
\end{tabular}
\end{center}

\index{possibility}
\texttt{EF} is ``possibly'', read forward in time: something can still happen.
Its negation is \texttt{AG not}: \texttt{never p} and \texttt{not EF p} are the
same question, so ``is it possible?'' and ``can it never happen?'' are one
question asked from both sides. A counterexample to a \texttt{never} is the
witness to the \texttt{EF}.

The atoms are what an observer can see: \texttt{\{"after": "Label"\}} holds in
the state a \texttt{Label} message or event has just reached, and
\texttt{"terminal"} holds where the session has ended. There are no atoms over a
role's internal state, because a protocol has none (Chapter~\ref{ch:language}).

\section{Asking \texttt{PlaceOrder}}

Four questions and LeanFM's answers, each with the formula as it is sent:

\begin{lstlisting}[language={},basicstyle=\ttfamily\footnotesize,breaklines=true]
Every session ends.
  {"AF":"terminal"}
  proved

A receipt is possible.
  {"EF":{"after":"Receipt"}}
  proved: witness: CreateOrder -> Quote -> Pay -> Charge -> Charged -> Receipt

Never charge after sold out.
  {"AG":{"implies":[{"after":"SoldOut"},{"AG":{"not":{"after":"Charge"}}}]}}
  proved

Every charge ends in a receipt.
  {"AG":{"implies":[{"after":"Charge"},{"AF":{"after":"Receipt"}}]}}
  disproved: counterexample: CreateOrder -> Quote -> Pay -> Charge -> Declined -> PaymentDeclined (and stays there)
\end{lstlisting}


The last one fails, and the counterexample says why: a charge can be declined.
The rule the author meant was probably ``every charge ends in a receipt or a
declined payment'', and the counterexample is what makes that visible.

\section{What an answer costs}

\index{model checking!cost}
Once the state graph exists, each CTL operator is one pass over it: reachability
for \texttt{EF}, strongly connected components for \texttt{EG} and \texttt{AF},
a fixpoint for \texttt{U}. The cost is about the number of states and transitions
times the size of the formula, and breadth-first search returns the shortest
witness, which is why the paths read well.

The size of the graph is the real cost. Concurrent branches interleave, a
deadline turns each state into a pair of state and elapsed time, and an attacker
choosing what to send multiplies every state by its choices. Two steps by
different roles that do not interact commute, so interleavings that differ only
by swapping them are the same behaviour; the explosion is real only where roles
share something. Every search has a bound, and an answer at the bound is
\texttt{resource-limit}, never \texttt{proved}. The server's \texttt{/metrics}
reports each check's time and model size, so the cost is measured, not guessed.

\section{Who can know what}

\index{knows}
\index{secret}
A security question asks what an attacker can derive. \texttt{knows(actor,
term)} decides whether an actor can build a term from everything it has seen,
taking pairs apart and opening ciphertexts with keys it can build, but never
guessing (the Dolev--Yao model). It is a state atom like \texttt{after}, so
\texttt{never knows(Ivo, nb)} is an ordinary CTL question.

Most of what an attacker learns does not matter, so a declaration marks what
must stay secret: a name ending in \texttt{\_secr}, or a term listed in
\texttt{secrets}. Each mark is the question \texttt{never knows(attacker, s)}.
Chapter~\ref{ch:csp-exercises} works through Needham--Schroeder, Lowe's attack,
and a password login whose hash leaks when it is sent in clear.

\section{Will two versions work together?}

\index{compatibility}
When a running protocol changes, not every role can be redeployed at once. For
each role LeanFM explores the mixed system, the new version of that role among
the old versions of the others, to every reachable state. Adding a branch in
which the shop answers \texttt{Backorder}:

\begin{lstlisting}[language={},basicstyle=\ttfamily\footnotesize,breaklines=true]
PlaceOrder → PlaceOrder: upgrade first: Client, Bank
Client cannot stay on the old version: Client waits to receive Quote from Shop, receive SoldOut from Shop; Shop is stuck trying to send Backorder to Client (after Client → Shop: CreateOrder; Shop chooses to send Backorder to Client)
Shop cannot upgrade first: Client waits to receive Quote from Shop, receive SoldOut from Shop; Shop is stuck trying to send Backorder to Client (after Client → Shop: CreateOrder; Shop chooses to send Backorder to Client)
\end{lstlisting}


The client and the bank can upgrade first; the shop cannot, because an old client
waits for a quote or a sold-out notice and would never accept a backorder.
